\begin{figure}[!ht]
% \vspace{-10pt}
    \centering
    \includegraphics[width=0.9\linewidth]{sec/figures/figure1.pdf}
    \caption{Illustration of how \ours\ evaluation works. Given the answers generated by an LVLM,
in step 1, we identify the descriptive content (with an LLM); 
In step 2, we extract corresponding atomic facts from the identified sentences; 
In step 3, the faithfulness of all atomic facts is verified according to the input image. 
In this example,  
the \underline{underlined part} denotes objective descriptive content in the answer.
The \textcolor{blue}{blue contents} denote hallucinations in the answers.
\ours\ allows a more
fine-grained and interpretable evaluation of the factual precision of free-form answers.
% , the result indicates that the response generated by LVLM-1 is better than that by LVLM-2.
}
\vspace{-5mm}
    \label{fig:hallucination}
\end{figure}


% \begin{window}[1,r,{\includegraphics[totalheight = 7cm]{sec/figures/figure1.pdf}},{Illustration of how \ours\ evaluation works. Given the answers generated by an LVLM,
% in step 1, we identify the descriptive content (with an LLM); 
% In step 2, we extract corresponding atomic facts from the identified sentences; 
% In step 3, the faithfulness of all atomic facts is verified according to the input image. 
% In this example,  
% The \underline{underlined part} denotes objective descriptive content in the answer.
% The \textcolor{blue}{blue contents} denote hallucinations in the answers.
% \ours\ allows a more
% fine-grained and interpretable evaluation of the factual precision of free-form answers.}]

% \end{window}



\section{Introduction}

% \begin{wrapfigure}{R}{0.5\textwidth} % This line sets the alignment to right and the width to half of the text width
%   \begin{center}
%   \includegraphics[width=0.48\textwidth]{sec/figures/figure1.pdf} % Adjust the width as needed
%     \end{center}
%   \caption{Illustration of how \ours\ evaluation works. Given the answers generated by an LVLM,
% in step 1, we identify the descriptive content (with an LLM); 
% In step 2, we extract corresponding atomic facts from the identified sentences; 
% In step 3, the faithfulness of all atomic facts is verified according to the input image. 
% In this example,  
% The \underline{underlined part} denotes objective descriptive content in the answer.
% The \textcolor{blue}{blue contents} denote hallucinations in the answers.
% \ours\ allows a more
% fine-grained and interpretable evaluation of the factual precision of free-form answers.}
%   \label{fig:hallucination}
% \end{wrapfigure}

Large Language Models (LLMs), such as GPT-3~\cite{gpt3} and ChatGPT~\cite{chatgpt}, have demonstrated various language modeling capabilities. Despite their achievements, they still lack the capacity to handle multimodal inputs effectively. As a result, a significant amount of research has shifted its focus towards Large Vision-Language Models (LVLMs)~\cite{llava,mPLUG-Owl,llava_rlhf}
by incorporating powerful LLMs~\cite{llama,vicuna} and Vision Foundation Models~(VFMs)~\cite{vit,DBLP:journals/corr/abs-2108-07258}. 
% \xd{better use VLMs throughout since we use VLM is the title}
% which enhance the Vision Foundation Models~(VFMs)~\cite{vit,DBLP:journals/corr/abs-2108-07258} by incorporating powerful LLMs~\cite{llama,vicuna}.
LVLMs have shown strong performance on various multimodal tasks, such as Visual Question Answering~\cite{vqatask}, Image Captioning~\cite{mscoco}, and Multimodal Conversation~\cite{llava}.  



Despite the effectiveness of LVLMs, the problem of hallucination is pervasive, often leading these models to generate fabricated information that is incongruent with the provided visual input~\cite{DBLP:conf/emnlp/RohrbachHBDS18,liu2023aligning,liu2023hallusionbench,DBLP:journals/corr/abs-2310-16045,chang2024unifiedhallucinationmitigationframework}. 
In the context of LVLM, the problem of hallucination can manifest as answers containing descriptions of the input image that are incorrect~\cite{Li2023EvaluatingOH}. 
As shown in Figure \ref{fig:hallucination}, the LVLM-2 generates an answer with an inaccurate description~(\ie \textit{standing on the front of a car}), which is not faithful towards the input image. The phenomenon of hallucination in LVLMs introduces potential hazards that could result in significant consequences such as misinformation and safety concerns, thus degrading the model's reliability in practical applications inevitably~\cite{DBLP:conf/chi/MacLeodBMC17}. Hence, it is imperative that these issues are thoroughly measured and addressed~\cite{ji2023survey}. 
% 
% \todo{motivation: the importance of the problem}

% To learn the hallucination level of different VLMs, limited work begins to evaluate VLMs with the new metric or human-annotated dataset in terms of hallucination. \citet{DBLP:journals/corr/abs-2305-10355} firstly proposed to evaluate hallucinations of VLMs. In particular, this work utilized caption hallucination assessment with image relevance (CHAIR)\citet{DBLP:conf/emnlp/RohrbachHBDS18} and Polling-based Object
% Probing Evaluation~(POPE) to evaluate object-level hallucination. Unlike previous work that
% only consider object hallucination, \citet{gunjal2023detecting} additionally manually annotate a dataset that contains both unfaithful entity descriptions and relationships. 
% In addition, \cite{DBLP:journals/corr/abs-2308-15126} proposed a hallucination evaluation framework
% based on LLM and human-annonated dataset. However, existing work either explores only object-level hallucination problems or only the image captioning task and is unable to evaluate a wide variety of free-form answers. 


% To delve into the hallucination levels across different LVLMs, 
% Despite the hallucinations in LVLMs should not be underestimated,
% Despite of the severe, 
Nevertheless, there have been limited explorations that 
% that employs new metrics or human-annotated datasets to 
measure the hallucination problem in LVLMs. \citet{Li2023EvaluatingOH} was among the first to measure the hallucinations of LVLMs with a polling-based object probing evaluation method. 
In addition, \citet{gunjal2023detecting} annotated a multi-modal hallucination detection dataset tailored for detailed image description evaluation. \citet{lovenia2023negative} devised Negative Object Presence Evaluation (NOPE), which used VQA to quantitatively evaluate object hallucination in LVLMs. 
% \citet{gunjal2023detecting} proposed to detect hallucinations in more detailed image captions and investigated utilizing a reward model for mitigating them. 
% \citet{DBLP:journals/corr/abs-2308-15126} introduced a hallucination evaluation framework grounded in LLMs and human-annotated datasets. 
% Although prior work has achieved success, 
These approaches, however, exhibit two key weaknesses:
(1) they focus on the limited setting of image captioning, and none of them explored evaluating hallucination of the complex and free-form answers to the open-ended questions~\cite{OpenAI2023GPT4V} (e.g. multimodal conversations~\cite{llava,DBLP:conf/acl-convai/SundarH22}, world knowledge-based VQA~\cite{AOKVQA} and visual storytelling~\cite{Huang2016VisualStory});
(2) they ignore fine-grained hallucinations of visual attributes in the generated answer \cite{liu2024mitigating,DBLP:conf/emnlp/RohrbachHBDS18}. 
% Yet, much of the existing research either narrowly addresses object-level hallucination or limits itself to the image captioning task, thus hindering a comprehensive evaluation of a broader range of free-form responses, such as complex questions~\cite{llava}.\xd{this paragraph is too long, summarize Rohrbah, Gunjal and Wang into one 3-4 sentences and call it prior work~\cite{DBLP:conf/emnlp/RohrbachHBDS18, gunjal2023detecting, DBLP:journals/corr/abs-2308-15126}}
% Another major limitation of work above is that they all tackle the setting of captioning (?), which is far from the distribution of real broad real user open-ended questions~\cite{OpenAI2023GPT4V}, whose answers are of free-form and more complex.

% \begin{figure}
%     \centering
%     \includegraphics[width=1.0\linewidth]{sec/figures/fig2.pdf}
%     \caption{A testing sample of complex reasoning question for LVLMs. The \textcolor{blue}{blue words} denotes commonsense reasoning-related content in the answer.}
%     \label{fig:commonsense}
% \end{figure}

% However, evaluating hallucination involved in free-form responses is non-trivial due to two reasons. (1) \textbf{Hybrid of Description and Commonsense.} Different from image captioning, the free-form response to the complex question not only generates descriptive content of the input image but also involves some common sense reasoning-related content. As shown in Figure~\ref{fig:commonsense}, the content of some sub-sentences (\ie sub-sentences in the blue font) does not need to be verified because it isn't the faithful description of the visual modality. Hence, how to identify the descriptive content in the response generated by VLMs is a key challenge. 
% (2) \textbf{Multiplicity of Hallucinations.} The existing methodology narrowly gauges hallucination, focusing on limited objects. It neglects other significant aspects such as counts, colors, and spatial relationships between objects, \etc Therefore, how to comprehensively evaluate various hallucinations is another significant challenge. 

Evaluating hallucinations present in free-form answer is especially challenging for two primary reasons:
(1) \textbf{Free-form answers contain a hybrid of descriptive and analytical contents.} Unlike close-domain tasks such as image captioning, answering open-domain questions in a free-form manner does not only require generating the question-relevant descriptive content of the given image. It may also involve analytical content such as rationales that include external commonsense knowledge. As depicted in Figure~\ref{fig:hallucination}, certain sub-sentences (\ie those without the underline) do not require verification with the image input due to their analytical nature. 
Because they encompass subjective analytical content that extends beyond a direct description of the visual inputs.
% they encompass subjective analytical content that extends beyond the visual inputs, rather than solely offering a direct description of the visual modality. 
Neglecting to distinguish between analytical and descriptive content inevitably distracts the factual measurement.
% Specifically, they contain commonsense reasoning content and do not offer a faithful description of the visual modality.
Thus, pinpointing the descriptive content within the answers generated by LVLMs is significant. 
% poses a major challenge. \xd{you mentioned the importance but not the challenge of identification, maybe remove this sentence or say it's ``important''}
(2) \textbf{Model outputs are prone to the multiplicity of hallucinations.} Current methodologies offer a constricted view on evaluating hallucinations, primarily concentrating on coarse-grained object existences~\cite{DBLP:conf/emnlp/RohrbachHBDS18,lovenia2023negative}, while neglecting other fine-grained elements, such as counts, colors, and the interrelations between objects (\eg the spatial relation between the person and the car in Figure~\ref{fig:hallucination}), which also form a significant portion of visual hallucinations~\cite{gunjal2023detecting}.
% For example, the relation information between the person and car in Figure~\ref{fig:hallucination}. 
Consequently, devising a method to holistically evaluate fine-grained hallucinations of visual attributes is also important.
% emerges as another substantial challenge\xd{similar here}.

% \todo{previous work focuses on ...}. 

% \begin{figure}
%     \centering
%     \includegraphics[width=\linewidth]{sec/figures/overview.png}
%     \caption{Illustration of how \ours\ works. Given the text answer generated by a VLM,
% \ours\ uses LLMs to identify the descriptive content and generate several atomic facts. Then faithfulness of all atomic facts is verified according to the input image. In this example, \ours\ allows a more
% fine-grained evaluation of factual precision, the result indicates that the answer generated by VLM 1 is better than that by VLM 2. The \textcolor{blue}{blue words} denotes commonsense reasoning-related content in the answer. \xd{merge with Figure 1}} 
%     \label{fig:overview}
% \end{figure}

% \todo{we propose \ours, ...}

% To overcome the above challenges, we propose \ours, which mainly consists of three components: descriptive sun-sentence identification, atomic fact generation, and fact verification, as shown in Figure~\ref{fig:ours}. The first component works on recognizing the descriptive sub-sentence from the hybrid content of the answer. The second component breaks the descriptive content into fine-grained elements (\ie facts), which contains various type of content, including color, relation, and attributes, \etc As for the last component, it focuses on verifying the consistency between the visual modality and atomic fact.

To address the aforementioned challenges, we propose the \ours\ metric, which can evaluate \textit{fine-grained hallucinations in various multimodal tasks}, such as image captioning and open-ended questions. This metric comprises three primary components: Descriptive Sub-sentence Identification, Atomic Fact Generation, and Fact Verification, as illustrated in Figure~\ref{fig:ours}. The first component is tasked with discerning descriptive sub-sentences within the composite content of the generated free-form answer. Thereafter, the second component deconstructs this descriptive content into fine-grained elements (\ie atomic facts)~\cite{factscore}. These atomic facts cover comprehensive types, such as objects attributes and interrelationships. The last component emphasizes verifying the consistency between the visual information and the derived atomic facts via a Visual Entailment Model (VEM)~\cite{DBLP:journals/corr/abs-1901-06706}. Based on the proposed metric, we evaluated several advanced LVLMs, such as LLaVA~\cite{llava} and MiniGPT-4~\cite{minigpt4}. From the results, we conclude that current LVLMs still face challenges of answers that are not faithful to the input image, which leaves a large room for improvement.

In summary, our contributions are as follows: 
(1) We introduce \ours, a metric tailored to assess hallucinations in LVLMs free-form answers to open-ended questions, which is not yet addressed by current studies; 
(2) To the best of our knowledge, we are the first to systematically study the LVLMs free-form answers and evaluate the fine-grained hallucinations of various types in LVLMs; 
(3) In our quest to understand the hallucinations manifested by LVLMs, we embark on comprehensive experiments with six open source models across diverse tasks and datasets. Our findings underscore that the hallucination phenomenon remains a pressing challenge for current LVLMs. As a byproduct, we released our code\footnote{\url{https://github.com/bcdnlp/FAITHSCORE}.}.
% \xd{too many usage of LLVMs in this paragraph so I replaced one with models}
% \begin{itemize}
%     \item We introduce \ours, a model tailored to assess various types of hallucinations in free-form answers, aspects often neglected by current methods.
%     \item To the best of our knowledge, this work is
% the first study that systematically evaluates the hallucination of free-form and open-domain answers in existing LVLMs. 
%     \item In our quest to understand the hallucinations manifested by LVLMs, we embarked on comprehensive experiments utilizing six publicly available VLMs across diverse datasets. Our findings underscore that addressing hallucination remains a pressing challenge for VLMs.
% \end{itemize}

